% 独立内容单元；对应原模块 4.1--4.3、4.6。
% 仅依赖主文档提供的标准数学、表格与 TikZ 宏包。
\section{灰度变换与直方图 从一个像素学到整幅图}

本节解决一个很具体的问题：照片里的东西明明存在，为什么看不清，怎样让它更容易看清？我们先学会给单个像素换一个数，再学会统计整幅图中的数，最后根据统计结果自动换数。对应原讲义的模块 4.1--4.3 和 4.6；原题 B1--B3、C1--C2、E1、F1、G3 都在本节保留题号并逐步解释。

\subsection{增强与复原 先分清想解决什么问题}

\textbf{图像增强（image enhancement）}是为了某个用途，改变图像的显示或特征。例如把较暗的车牌提亮，使字符容易辨认；把边缘突出，使零件轮廓容易测量。这里的“好”与用途有关：照片看起来自然，和机器容易读出字符，可能需要不同的处理。

\textbf{图像复原（image restoration）}是根据图像如何被破坏的模型，估计未被破坏时的图像。例如知道相机在曝光时水平晃动，就尝试根据运动模糊模型减轻拖影；知道噪声具有某种统计规律，就据此估计较干净的图像。

\textbf{先验知识（prior knowledge）}就是处理前已经掌握的信息，如运动方向、噪声特点或镜头模糊程度。“复原使用模型”不代表一定能恢复真实原图。数据可能已经丢失，模型也可能不准确。增强也可以有客观评价指标，例如字符识别正确率，因此不能把二者简单分成“完全主观”和“完全客观”。

想象同一张模糊的黑白照片：增大文字与背景的灰度差，是增强；估计造成拖影的过程再尝试反向补偿，是复原。两者也可以串联使用。

\subsubsection*{原题 F1 图像增强与图像复原有什么区别}
\label{original:F1}

\textbf{第一步，比较目标。}增强问“怎样让它更适合当前用途”；复原问“怎样从退化结果估计原来的图像”。

\textbf{第二步，比较依据。}增强可以根据观察效果、应用指标来选择变换；复原通常依赖退化模型、噪声模型等先验知识。

\textbf{第三步，举出例子。}低对比度照片的灰度拉伸属于增强；根据已知运动方向和程度进行去模糊属于复原。去噪可能属于增强流程，也可能建立在复原模型上，不能仅凭“去噪”这个名字判断。

\textbf{可写在答卷上的答案：}图像增强以特定应用为目标，改善可视效果或后续分析效果；图像复原依据退化过程的先验或模型，估计接近未退化状态的图像。二者的目标和依据不同，但可能使用相似工具，也都不保证恢复已经丢失的信息。

\textbf{处理途径也要分清。}直接在像素网格上处理称为\textbf{空间域处理（spatial-domain processing）}：只根据当前像素灰度换数，叫点处理；结合周围一小片像素计算，属于邻域或掩模处理，例如平滑与锐化。另一条途径是\textbf{频率域处理（frequency-domain processing）}：先把图像变换为不同变化频率的分量，再决定哪些保留、哪些减弱，最后变回图像。后面的滤波章节会解释后两类；本节先把点处理学明白。

\subsection{点运算的共同语言 输入一个数 输出一个数}

\textbf{灰度（gray level）}在本节中是表示像素明暗的数。对常见的 8 位无符号灰度图，0 表示黑，255 表示白，中间的数表示不同程度的灰。\textbf{灰度级数}记作 $L$，所以 8 位图有 $L=256$ 级，最大灰度为 $L-1=255$。注意“256 级”包含 0，因此最大数不是 256。

\textbf{点运算（point operation）}是根据某个像素自己的灰度，计算它的新灰度。把输入灰度记为 $r$，输出灰度记为 $s$，变换规则记为 $T$，就得到
\[
s=T(r).
\]
括号的意思是“把 $r$ 交给规则 $T$ 处理”，不是乘法。若规则是 $T(r)=r+20$，输入 $r=40$，便得到 $s=60$。像素仍在原位置，只是灰度变了。

画变换曲线时，横轴放输入 $r$，纵轴放输出 $s$。找一个输入，竖直向上碰到曲线，再水平向左读纵轴，就是输出。曲线在 $s=r$ 的对角线上方，该输入就变亮；在下方，该输入就变暗。

\textbf{全局点变换（global point transform）}使用全图共用的一条规则：两个位置的输入都等于 40，输出就相同。\textbf{局部处理（local processing）}会利用像素周围的信息；即使输入相同，不同位置也可能得到不同输出。HE 的规则由整幅图统计得到，但规则一旦确定，对所有像素使用同一张灰度对照表。

\subsubsection*{计算结果超出范围时怎么办}

计算可以产生小数，也可以得到负数或超过 255 的数。若最终保存为 8 位图，通常先\textbf{截断到合法范围（clipping）}，再\textbf{取整（rounding）}。例如 $260\to255$，$-7\to0$，$52.7\to53$。本节的 $\operatorname{round}(x)$ 表示把 $x$ 四舍五入到最近整数；题目若给出其他取整规则，要按题目执行。

截断会丢失差别：原本算得 260 与 300 的两个像素都变成 255。\textbf{饱和（saturation）}指许多不同的真实亮度被记录成同一个最大值或最小值。若一片白色区域的像素已经全部是 255，仅靠灰度变换无法知道它原来有哪些不同细节。

\subsection{逐个理解常见灰度变换}

\subsubsection*{线性变换 明暗平移与对比度缩放}

图像处理中常把 $s=ar+b$ 称为\textbf{线性灰度变换（linear gray-level transform）}；严格数学名称是\textbf{仿射变换（affine transform）}，当 $b\ne0$ 时不满足线性算子的定义。其中 $a$ 是乘在输入上的系数，$b$ 是最后加上的偏移量。先乘 $a$，再加 $b$，最后检查范围。

若 $a=1,b=30$，灰度 40 和 50 变成 70 和 80。两者都变亮，但差值仍是 10。若 $a=2,b=0$，它们变成 80 和 100，差值从 10 变成 20。第二种操作拉大了这两个灰度的区别。\textbf{对比度（contrast）}与明暗差异有关，不等于平均亮度。

这里的结论以没有截断为前提。把 140 和 160 都乘 2 后保存为 8 位，两个结果都截断成 255，差异反而消失。因此“系数越大，对比度越好”不成立。

\subsubsection*{反色 把灰度次序翻过来}

\textbf{图像反色或负片变换（image negative）}把黑变白、白变黑，中间灰度反向对应。公式是
\[
s=(L-1)-r.
\]
$L-1$ 是最大灰度，减号表示从最大值中减去输入。例如 8 位图中，$r=30$ 时 $s=255-30=225$；$r=200$ 时 $s=55$。输入越大，输出越小，所以曲线是一条下降直线。

反色改变黑白呈现方式，有时能让原先不容易观察的细节更醒目。但在未取整的情况下，两灰度的差值大小不变，仅次序颠倒；它没有凭空增加细节，也不等于一般意义上的“提高对比度”。

\subsubsection*{对数变换 压缩跨度很大的数}

\textbf{对数变换（logarithmic transform）}让输入越大时，输出增长得越慢。先不记公式，观察这个规律：从 1 增加到 10 是乘 10，从 10 增加到 100 也是乘 10；以 10 为底的对数把它们变成 $0,1,2$，乘法尺度被压成相近的加法尺度。

本节用自然对数 $\ln$，其底数是常数 $e\approx2.718$。常见灰度公式为 $s=c\ln(1+r)$：先加 1，再取对数，最后乘缩放系数 $c$。加 1 是为了让黑色 $r=0$ 可以安全计算，且 $\ln(1)=0$。

若输入范围为 $[0,L-1]$，又希望最大输入仍映射到最大输出，就令
\[
c=\frac{L-1}{\ln L},\qquad
s=\frac{L-1}{\ln L}\ln(1+r).
\]
方括号 $[a,b]$ 表示从 $a$ 到 $b$、含两端的范围。把 $r=L-1$ 代入，$1+r=L$，于是 $\ln L$ 与分母相除为 1，输出确实是 $L-1$。

用只有 8 级的图说明：$L=8$，所以 $c=7/\ln8\approx3.366$。输入 1 对应 $3.366\ln2\approx2.333$，输入 3 对应 $3.366\ln4\approx4.667$，输入 7 对应 7。较暗的输入被明显抬高，而亮部数值被挤得更近。

常见应用是显示\textbf{动态范围（dynamic range）}很大的数据，如傅里叶频谱幅度。若大部分数较小，少数数极大，直接显示可能只看见几个亮点；对数压缩后，较小数值也更容易观察。这里改善的是显示，不是证明微弱细节都是真实信号。

\subsubsection*{幂律与伽马变换 先归一化再判断明暗}

\textbf{幂律变换（power-law transform）}使用指数来弯曲灰度曲线；在图像处理中常称\textbf{伽马变换（gamma transform）}。指数记为希腊字母 $\gamma$，读作 gamma。例如 $u^2$ 表示 $u\times u$，$u^{1/2}$ 表示 $u$ 的平方根。

为了避免单位和范围造成误判，先把输入除以最大灰度：$u=r/(L-1)$，使 $u$ 落在 0 到 1。这个步骤叫\textbf{归一化（normalization）}。然后算 $v=u^\gamma$，最后乘回最大灰度：
\[
s=(L-1)\left(\frac{r}{L-1}\right)^\gamma,\qquad \gamma>0.
\]
$r$ 是原灰度，$u$ 是归一化输入，$v$ 是归一化输出，$s=(L-1)v$。式子从内到外读：先做除法，再做乘方，最后做乘法。

看三个容易手算的归一化输入：
\begin{center}
\begin{tabular}{ccccc}
\toprule
输入 $u$ & $u^{1/2}$ & $u^1$ & $u^2$ & 判断 \\
\midrule
$0.25$ & $0.50$ & $0.25$ & $0.0625$ & 平方根提亮，平方压暗 \\
$0.50$ & $0.7071$ & $0.50$ & $0.25$ & 规律相同 \\
$1$ & $1$ & $1$ & $1$ & 白端保持不变 \\
\bottomrule
\end{tabular}
\end{center}
例如 8 位图 $r=64,\gamma=1/2$：先算 $64/255\approx0.251$，取平方根约为 $0.501$，乘 255 得约 $127.75$，取整为 128。

所以在上述规范化公式下，$0<\gamma<1$ 提亮中间灰度，$\gamma=1$ 不变，$\gamma>1$ 压暗中间灰度。黑端和白端保持不变。\textbf{不能直接把 $s=r^\gamma$ 与这一结论混用。}例如直接对原始灰度算 $64^{1/2}=8$，当然不是提亮；它缺少了范围归一化与输出缩放。

有的资料写 $s=cr^\gamma$。若想保持两个端点，需要选 $c=(L-1)^{1-\gamma}$，这才与上面的规范化公式一致。另有软件把控制参数写成指数的倒数，因此使用软件时必须先看实际公式，不能只看“gamma 大于 1”这几个字。

伽马变换也会改变某些灰度区间的间距，所以可能改善局部灰度对比。但全图统一伽马使用同一条曲线，不能根据空间位置分别照顾暗角与亮区。\textbf{伽马校正（gamma correction）}还可用于补偿采集或显示设备的非线性响应；它与“把照片提亮”有关，但用途并不完全相同。

\subsubsection*{原题 C2 画反色 对数与两种幂律曲线}
\label{original:C2}

\textbf{先画基准。}以下使用归一化坐标 $u=r/(L-1)$、$v=s/(L-1)$，因此两轴均从 0 到 1。换回题目要求的原灰度，只要把轴端 1 改标为 $L-1$，形状不变。虚线 $v=u$ 是不作改变的参照线。

\begin{center}
\begin{tikzpicture}[x=3.9cm,y=3.9cm,>=Stealth,font=\small]
\begin{scope}
\draw[->] (0,0)--(1.12,0) node[right] {$u$};
\draw[->] (0,0)--(0,1.12) node[above] {$v$};
\draw[dashed] (0,0)--(1,1);
\draw[very thick] (0,1)--(1,0);
\node[below left] at (0,0) {0};\node[below] at (1,0) {1};\node[left] at (0,1) {1};
\node[align=center] at (.5,-.24) {反色\quad $v=1-u$};
\end{scope}
\begin{scope}[xshift=7.3cm]
\draw[->] (0,0)--(1.12,0) node[right] {$u$};
\draw[->] (0,0)--(0,1.12) node[above] {$v$};
\draw[dashed] (0,0)--(1,1);
\draw[very thick,domain=0:1,samples=60] plot (\x,{ln(1+7*\x)/ln(8)});
\node[below left] at (0,0) {0};\node[below] at (1,0) {1};\node[left] at (0,1) {1};
\node[align=center] at (.5,-.24) {对数示例\quad $v=\ln(1+7u)/\ln8$};
\end{scope}
\begin{scope}[yshift=-6.1cm]
\draw[->] (0,0)--(1.12,0) node[right] {$u$};
\draw[->] (0,0)--(0,1.12) node[above] {$v$};
\draw[dashed] (0,0)--(1,1);
\draw[very thick,domain=0:1,samples=70] plot (\x,{sqrt(\x)});
\node[below left] at (0,0) {0};\node[below] at (1,0) {1};\node[left] at (0,1) {1};
\node[align=center] at (.5,-.24) {幂律\quad $v=u^{1/2}$\quad 提亮};
\end{scope}
\begin{scope}[xshift=7.3cm,yshift=-6.1cm]
\draw[->] (0,0)--(1.12,0) node[right] {$u$};
\draw[->] (0,0)--(0,1.12) node[above] {$v$};
\draw[dashed] (0,0)--(1,1);
\draw[very thick,domain=0:1,samples=60] plot (\x,{\x*\x});
\node[below left] at (0,0) {0};\node[below] at (1,0) {1};\node[left] at (0,1) {1};
\node[align=center] at (.5,-.24) {幂律\quad $v=u^2$\quad 压暗};
\end{scope}
\end{tikzpicture}
\end{center}

\textbf{逐条作图：}反色连结左上角与右下角；对数从原点到右上角，前段较陡、后段较平；$\gamma<1$ 的规范化幂律位于对角线上方；$\gamma>1$ 的位于下方。用“上方/下方、越来越陡/越来越平”描述最不容易混淆。若使用数学凹凸术语，对数与 $0<\gamma<1$ 的幂律在内部是凹函数，$\gamma>1$ 的幂律是凸函数。

\subsubsection*{分段线性变换 给不同区间不同的待遇}

\textbf{分段线性变换（piecewise-linear transform）}把输入范围切成几段，每段使用一条直线。\textbf{对比度拉伸（contrast stretching）}是其中一个常见用途：把原来挤在较窄范围里的灰度拉开。

设希望输入 $a$ 映到黑色 0，输入 $b$ 映到白色 $L-1$，并且 $a<b$。中间灰度先减去 $a$，把起点移到 0；再除以 $b-a$，换成 0 到 1 之间的相对位置；最后乘 $L-1$：
\[
s=\begin{cases}
0,&r\le a,\\[2pt]
(L-1)\dfrac{r-a}{b-a},&a<r<b,\\[5pt]
L-1,&r\ge b.
\end{cases}
\]
大括号表示“按条件选一行”，不是三行全部相加。例：$L=256,a=80,b=160$，输入 120 位于两端正中，因此 $s=255(120-80)/(160-80)=127.5$，取整为 128。输入 100 变为约 64，输入 140 变为约 191，原来的 40 级间距变成约 127 级。

\begin{center}
\begin{tikzpicture}[x=.045cm,y=.018cm,>=Stealth,font=\small]
\draw[->] (0,0)--(275,0) node[right] {$r$};
\draw[->] (0,0)--(0,290) node[above] {$s$};
\draw[dashed] (0,0)--(255,255);
\draw[very thick] (0,0)--(80,0)--(160,255)--(255,255);
\draw[dotted] (160,0)--(160,255)--(0,255);
\node[below] at (80,0) {80};\node[below] at (160,0) {160};\node[below] at (255,0) {255};
\node[left] at (0,255) {255};\node[below left] at (0,0) {0};
\node[align=center,fill=white] at (211,130) {斜率较大的区段\\拉开灰度差};
\end{tikzpicture}
\end{center}

这条曲线牺牲了 $a$ 以下和 $b$ 以上的差别。若两端也有重要细节，可以设置多个控制点，如 $(0,0),(r_1,s_1),(r_2,s_2),(L-1,L-1)$，用直线连接，给中间重点区间较大斜率，给两端较小斜率。这里 $(r_1,s_1)$ 表示“输入 $r_1$ 对应输出 $s_1$”。

\subsubsection*{灰度级切片 突出一个数值区间}

\textbf{灰度级切片（gray-level slicing）}把落在指定区间中的像素突出显示。设目标范围为 $a\le r\le b$，想让这些像素显示为白色。有两种常见做法：
\begin{enumerate}[leftmargin=*]
\item \textbf{其余归零：}范围内输出 $L-1$，范围外输出 0。结果只有黑白，便于显示目标分布。
\item \textbf{保留背景：}范围内输出 $L-1$，范围外仍输出原值 $r$。目标突出，同时保留环境背景。
\end{enumerate}
例如取区间 $[100,140]$，输入 $90,120,170$，第一种输出 $0,255,0$，第二种输出 $90,255,170$。选择的是\textbf{灰度范围}，不是自动识别“哪个物体”：如果衣服和墙面恰好具有相同灰度，它们会一起被选中。

\subsubsection*{位平面切片 拆开一个数的二进制组成}

\textbf{位平面切片（bit-plane slicing）}按二进制位拆分图像，与“挑选一个灰度区间”完全不同。\textbf{二进制（binary）}只有 0 和 1；各位的权重依次是 $1,2,4,8,16,32,64,128$。

例如十进制灰度 173 可以拆成 $128+32+8+4+1$，其 8 位二进制写为 $10101101$。从右到左把位编号为 $k=0,1,\ldots,7$，该位的值记为 $b_k$，只会等于 0 或 1。于是
\[
r=\sum_{k=0}^{7}b_k2^k
=b_0\!\times1+b_1\!\times2+\cdots+b_7\!\times128.
\]
符号 $\sum$ 表示“把后面这类项按指定编号逐项相加”；下方 $k=0$ 是起点，上方 7 是终点。$2^k$ 是第 $k$ 位的权重。

把全图每个像素的同一位取出来，就得到一个位平面；显示时通常把位值 0 显示为黑，把位值 1 显示为白。若想用四则运算提取第 $k$ 位：先除以 $2^k$，向下取整，再除以 2 取余数。例如提取 173 的第 5 位：$173/32=5.40625$，向下取整得 5，5 除以 2 的余数为 1。

只保留 173 的最高四位，相当于保留 $128+32=160$，舍掉了 13。高位通常决定较大的明暗结构，低位影响较细的数值变化；\textbf{低位不一定全是噪声}，也可能含有真实微弱细节。是否丢弃，要看任务允许多大误差。

\subsection{直方图 从数像素开始}

\textbf{灰度直方图（gray-level histogram）}回答“每种灰度有多少像素”。它不记录像素在哪里。横轴是灰度，纵轴是像素数，每个灰度对应的计数格也称\textbf{组或箱（bin）}。

先看一幅只有 10 个像素、允许灰度为 $0,1,2,3$ 的小图，其灰度依次为
\[
0,\ 0,\ 1,\ 1,\ 1,\ 1,\ 1,\ 2,\ 3,\ 3.
\]
数一数：0 出现 2 次，1 出现 5 次，2 出现 1 次，3 出现 2 次。因此直方图计数是 $2,5,1,2$。把灰度为 $k$ 的个数记为 $n_k$，把总像素数记为 $n$，本例有 $n=10$。

\textbf{归一化直方图（normalized histogram）}把每个计数除以总数。这个比值可以看作：从图中等可能地随机选一个像素，它恰好为该灰度的\textbf{概率（probability）}。无需先学概率论，记住“这一类有多少，除以总共有多少”即可：
\[
p_k=\frac{n_k}{n}.
\]
$p_k$ 是灰度 $k$ 所占的比例。例如灰度 1 的比例是 $5/10=0.5=50\%$。所有比例相加必须等于 1，因为每个像素恰好属于一个灰度。

\textbf{累积分布函数（cumulative distribution function，CDF）}进一步回答“灰度不大于 $k$ 的像素，占全部的多少”。所谓累计，就是从最左边一直加到当前位置。把累计比例记为 $C_k$，累计个数记为 $N_k$：
\[
N_k=n_0+n_1+\cdots+n_k,\qquad
C_k=\frac{N_k}{n}=p_0+p_1+\cdots+p_k.
\]
例如“不大于 1”包含灰度 0 和 1，两类共 $2+5=7$ 个，所以 $C_1=7/10=0.7$，不是 $5/10$。

\begin{center}
\begin{tabular}{ccccc}
\toprule
灰度 $k$ & 计数 $n_k$ & 概率 $p_k$ & 累计个数 $N_k$ & 累计比例 $C_k$\\
\midrule
0 & 2 & $2/10=0.2$ & 2 & $2/10=0.2$\\
1 & 5 & $5/10=0.5$ & $2+5=7$ & $7/10=0.7$\\
2 & 1 & $1/10=0.1$ & $7+1=8$ & $8/10=0.8$\\
3 & 2 & $2/10=0.2$ & $8+2=10$ & $10/10=1$\\
\bottomrule
\end{tabular}
\end{center}

\begin{center}
\begin{tikzpicture}[x=1.1cm,y=.55cm,>=Stealth,font=\small]
\begin{scope}
\draw[->] (-.5,0)--(4,0) node[right] {$k$};
\draw[->] (-.5,0)--(-.5,6.2) node[above] {计数 $n_k$};
\foreach \k/\v in {0/2,1/5,2/1,3/2}{
\draw[fill=black!15] (\k-.28,0) rectangle (\k+.28,\v);
\node[above] at (\k,\v) {\v};\node[below] at (\k,0) {\k};}
\end{scope}
\begin{scope}[xshift=8cm,y=3cm]
\draw[->] (-.5,0)--(4,0) node[right] {$k$};
\draw[->] (-.5,0)--(-.5,1.15) node[above] {累计比例 $C_k$};
\draw[dotted] (-.5,1)--(3.5,1);\node[left] at (-.5,1) {1};
\draw[very thick] (0,.2)--(1,.2)--(1,.7)--(2,.7)--(2,.8)--(3,.8)--(3,1)--(3.5,1);
\foreach \k/\v in {0/.2,1/.7,2/.8,3/1}{\fill (\k,\v) circle (1.7pt);\node[above left] at (\k,\v) {\v};\node[below] at (\k,0) {\k};}
\end{scope}
\end{tikzpicture}
\end{center}

左图每根柱子单独计数；右图不断把新柱子的比例加进来。因此 CDF 只能上升或保持不变，不能下降，最右端必为 1。某灰度没有像素，该处累计值就不增加。\textbf{高柱子对应累计曲线较大的跳跃}，这是后面均衡化会放大某些小灰度差的原因。

\subsubsection*{直方图能告诉什么 不能告诉什么}

直方图偏左，说明低灰度像素占比大，图像整体常显得偏暗；偏右则常偏亮。集中在窄范围，通常说明可利用的灰度跨度较小；分布很宽，说明存在跨度较大的灰度。但少数黑点和白点就能把范围撑满，不能据此断言每个局部都清楚。

如果把同一批像素打乱位置，直方图不会改变，但一张人像可以被打乱成杂乱斑点。因此直方图\textbf{没有空间结构信息}，不能单独判断清晰度、目标位置、噪声种类，也不能证明图中恰好只有两个物体。

\subsection{直方图均衡 为什么累计比例能变成新灰度}

\textbf{直方图均衡化（histogram equalization，HE）}用输入灰度的累计比例来决定输出灰度。直觉上，它把“我排在全部像素的什么位置”变成“我应该落在输出灰度范围的什么位置”。

假设累计比例为 0.7，表示有 70\% 的像素不比当前灰度亮；让这个灰度落在输出范围约 70\% 的位置，就是均衡映射的基本想法。对于 $L$ 个输出灰度级，本节采用的离散公式为
\[
s_k=\operatorname{round}\!\bigl((L-1)C_k\bigr).
\]
$s_k$ 表示原灰度 $k$ 应换成的新灰度。$C_k$ 是 0 到 1 的累计比例，乘 $L-1$ 把它扩大到合法灰度范围，再取整得到可保存的整数。

刚才的 10 像素例子有 $L=4$，因此乘数为 3：
\begin{center}
\begin{tabular}{ccccc}
\toprule
输入 $k$ & $C_k$ & $3C_k$ & 取整输出 $s_k$ & 这些像素搬到哪里\\
\midrule
0 & 0.2 & 0.6 & 1 & 2 个全部移到 1\\
1 & 0.7 & 2.1 & 2 & 5 个全部移到 2\\
2 & 0.8 & 2.4 & 2 & 1 个也移到 2\\
3 & 1 & 3 & 3 & 2 个保持在 3\\
\bottomrule
\end{tabular}
\end{center}
输出灰度 $0,1,2,3$ 的计数为 $0,2,6,2$。灰度 2 收到原灰度 1 的 5 个像素和原灰度 2 的 1 个像素，所以为 $5+1=6$。\textbf{均衡化不是把像素重新放到空间位置，而是把它们的灰度重新编号。}

为什么名字叫“均衡”，结果却不平均？在灰度连续、分布满足相应条件的理想模型中，用连续 CDF 变换会产生均匀分布；实际图像有有限个整数灰度，同值像素又被同一个规则整体映射，不能照搬这个理想结论。上面的例子反而出现了更高的柱子，正好提醒我们：\textbf{HE 不保证每根柱子更接近平均高度，也不保证视觉效果一定改善。}

\textbf{软件结果不同时先查公式。}有的实现还会减去第一个非零累计值，再缩放到全范围；取整规则也可能不同。本节原题 E1 明确采用 $\operatorname{round}((L-1)C_k)$，计算中不要擅自换成另一种版本。常数图像等特殊输入还需要单独约定，不能把所有实现都当成一模一样。

\subsubsection*{进阶阅读 连续公式中的积分到底在累计什么}

原讲义还给出连续公式。若暂时没有学过积分，可以先把它看作“柱子无限变窄时，累计柱子面积”的写法，不需要先会求复杂积分。

把连续输入记为 $r$，并先归一化到 $[0,1]$。\textbf{概率密度函数（probability density function，PDF）}记为 $p_R(r)$，是描述灰度分布的密度曲线。离散图像的 $p_k$ 是某一个灰度的概率；连续情形中，\textbf{一个区间下方的面积}才是落在该区间的概率，曲线高度本身不是单点概率。曲线总面积为 1。

从 0 累计到当前灰度 $r$ 的面积就是 CDF：
\[
s=T(r)=\int_0^r p_R(w)\,\mathrm{d}w.
\]
弯曲的符号 $\int$ 表示积分，也就是连续累计；下方 0 是开始位置，上方 $r$ 是结束位置；$w$ 是累计过程中走过的临时灰度；$p_R(w)$ 是那个位置的密度；$\mathrm{d}w$ 表示极小的一段灰度宽度。可把 $p_R(w)\mathrm{d}w$ 想成一根极细小柱的面积，再把所有小面积加起来。输出 $s$ 在 0 到 1 之间。

例如输入在 $[0,1]$ 上本来就均匀，密度为 1；从 0 到 $r$ 的面积是“高 1 乘宽 $r$”，于是 $T(r)=r$，无需改变。若满足连续分布条件，CDF 变换后的 $s$ 在 $[0,1]$ 上均匀，写成 $p_S(s)=1$；$p_S$ 表示输出的密度。若输出范围改为 $[0,L-1]$，需在 CDF 外乘 $L-1$，均匀密度也相应为 $1/(L-1)$，不能仍写成 1。

实际数字图像要把积分换成累计和，再量化成整数；这一步破坏了连续结论所需的条件，所以不能要求实际输出逐级等高。

\subsubsection*{原题 E1 完整计算 64 乘 64 图像的均衡映射}
\label{original:E1}

\textbf{题目。}一幅 $64\times64$ 的图像只有 8 个允许的灰度级 $0,1,\ldots,7$，其像素计数如下，求均衡后的映射关系与新直方图。
\begin{center}
\begin{tabular}{c|rrrrrrrr}
\toprule
灰度 & 0&1&2&3&4&5&6&7\\
计数 &790&1023&850&656&329&245&122&81\\
\bottomrule
\end{tabular}
\end{center}

\textbf{第 1 步，确认两个不同的数量。}总像素数 $n=64\times64=4096$，用于求比例；灰度级数 $L=8$，所以输出最大值 $L-1=7$，用于缩放。4096 是“有多少个像素”，8 是“每个像素允许取几种数”，不能互换。

先验算计数总和：$790+1023+850+656+329+245+122+81=4096$，与图像尺寸一致。

\textbf{第 2 步，从左到右求累计个数。}第 0 项为 790；第 1 项是 $790+1023=1813$；第 2 项再加 850 得 2663。每一步都只需在前一累计数上加当前计数。最后必须得到 4096。

\textbf{第 3 步，用精确计数算比例和输出。}比如输入灰度 2：累计个数 2663，累计比例为 $2663/4096\approx0.650146$，乘 7 得 $4.551025$，四舍五入后输出 5。表中的小数用于展示，真正决定取整的是精确分数 $7N_k/4096$。

\begin{center}
\small
\begin{tabular}{rrrrrrr}
\toprule
$k$ & $n_k$ & $N_k$ & $p_k=n_k/4096$ & $C_k=N_k/4096$ & $7C_k$ & $s_k$\\
\midrule
0&790&790&0.192871&0.192871&1.350098&1\\
1&1023&1813&0.249756&0.442627&3.098389&3\\
2&850&2663&0.207520&0.650146&4.551025&5\\
3&656&3319&0.160156&0.810303&5.672119&6\\
4&329&3648&0.080322&0.890625&6.234375&6\\
5&245&3893&0.059814&0.950439&6.653076&7\\
6&122&4015&0.029785&0.980225&6.861572&7\\
7&81&4096&0.019775&1.000000&7.000000&7\\
\bottomrule
\end{tabular}
\end{center}

不要先把概率都保留两位小数再累计。粗略四舍五入可能把某一结果推过半整数分界，改变最终映射。本题恰好得到与原讲义相同的整数映射，不代表提前舍入总是安全。

\textbf{第 4 步，写成灰度对照表。}
\[
0\to1,\quad 1\to3,\quad 2\to5,\quad
3\to6,\quad 4\to6,\quad
5\to7,\quad 6\to7,\quad 7\to7.
\]
箭头表示“每个输入灰度换成哪个输出灰度”。例如原图左上角若是灰度 4，就把该位置写成 6；不需要移动那个像素。

\textbf{第 5 步，按输出灰度合并计数。}输出 6 接收输入 3 和 4，所以计数为 $656+329=985$。输出 7 接收输入 5、6、7，所以计数为 $245+122+81=448$。没人映射到 0、2、4，它们的计数就是 0。

\begin{center}
\begin{tabular}{c|rrrrrrrr}
\toprule
输出灰度 &0&1&2&3&4&5&6&7\\
新计数 &0&790&0&1023&0&850&985&448\\
\bottomrule
\end{tabular}
\end{center}
核查：$790+1023+850+985+448=4096$，总像素没有增减。若总数不对，通常是合并时漏了一组，或把灰度值误当作像素数来加。

\begin{center}
\begin{tikzpicture}[x=.66cm,y=.0034cm,>=Stealth,font=\small]
\begin{scope}
\draw[->] (-.6,0)--(8,0) node[right] {灰度};
\draw[->] (-.6,0)--(-.6,1210) node[above] {计数};
\node[above] at (3.5,1200) {E1 输入直方图};
\foreach \k/\v in {0/790,1/1023,2/850,3/656,4/329,5/245,6/122,7/81}{\draw[fill=black!15] (\k-.32,0) rectangle (\k+.32,\v);\node[below] at (\k,0) {\k};}
\draw[dotted] (-.6,1023)--(7.5,1023);\node[left] at (-.6,1023) {1023};
\end{scope}
\begin{scope}[xshift=7.8cm]
\draw[->] (-.6,0)--(8,0) node[right] {灰度};
\draw[->] (-.6,0)--(-.6,1210) node[above] {计数};
\node[above] at (3.5,1200) {E1 输出直方图};
\foreach \k/\v in {0/0,1/790,2/0,3/1023,4/0,5/850,6/985,7/448}{\draw[fill=black!15] (\k-.32,0) rectangle (\k+.32,\v);\node[below] at (\k,0) {\k};}
\draw[dotted] (-.6,1023)--(7.5,1023);
\end{scope}
\end{tikzpicture}
\end{center}

\textbf{第 6 步，解释结果而非只报数字。}输出占用的灰度级从 8 个变为 5 个，某些级合并、某些级空缺。若严格均匀，每一级应有 $4096/8=512$ 个像素；实际有 0、790、1023 等，显然不均匀。输入原本就有像素覆盖 0 到 7，所以也不能声称这道题“把原有范围变得更宽”。它实现了规定的累计映射，至于是否更适合观看或识别，还要检查图像和任务。

\subsubsection*{原题 B3 为什么均衡后不是完全均匀}
\label{original:B3}

\textbf{第一步，区分理想模型与数字图像。}连续 CDF 均衡的结论有连续分布等条件；数字图像只有有限个灰度级，需取整。

\textbf{第二步，抓住同值像素必须一起变换。}E1 中 1023 个输入灰度 1 的像素都被映射成 3，不能由这张单值灰度表把它们拆成若干份分配到不同灰度。所以只要输入某个计数很大，该组就不会被平均拆散。

\textbf{第三步，说明取整造成合并。}输入 3 与 4 的未取整输出分别约为 5.672 和 6.234，取整后都变为 6，于是计数相加。出现空级和合并级是正常结果。

\textbf{完整结论：}标准离散 HE 使用确定性的灰度映射，像素同值成组搬移，且输出要量化到整数级，因此一般不会得到完全均匀的直方图。它也不保证按某个均匀性指标一定更接近均匀；不能把“名字叫均衡”当作结果保证。

\subsection{直方图匹配 让图像接近指定的灰度分布}

\textbf{直方图规定化或匹配（histogram specification or matching）}是在给定目标分布后，让输出的灰度分布尽量接近该目标。均衡化使用固定的累计映射思路；匹配允许目标偏暗、偏亮、单峰或其他指定形状。

先理解“相同累计位置”：若源灰度 1 的累计比例为 0.7，就去目标分布中找累计比例接近 0.7 的灰度。它们都处在各自分布大约 70\% 的位置，把两者对应起来即可。

源灰度记为 $r$，它的 CDF 记为 $T(r)$；目标灰度记为 $z$，目标 CDF 记为 $G(z)$。先求 $s=T(r)$，再寻找满足 $G(z)\approx s$ 的目标灰度。在可逆的理想连续情形，可简记为
\[
z=G^{-1}(T(r)).
\]
这里 $G^{-1}$ 表示“沿着函数反查输入”，\textbf{不是 $1/G$}。数字图像的 CDF 常有台阶和平台，通常没有普通意义上的唯一逆函数，需要明确近似匹配规则。

\textbf{一个可以手算的离散例子。}仍用前面 10 像素小图，源累计比例是 $0.2,0.7,0.8,1$。希望目标灰度 $0,1,2,3$ 的概率为 $0.1,0.2,0.3,0.4$，目标累计比例便为 $0.1,0.3,0.6,1$。

本例采用“选累计比例差的绝对值最小者，若并列则选较低灰度”。绝对值 $|a-b|$ 表示两个数相差多远，例如 $|0.7-0.6|=0.1$。
\begin{center}
\begin{tabular}{cccl}
\toprule
源灰度 & 源累计值 & 所选目标灰度 & 比较过程\\
\midrule
0&0.2&0&到 0.1 和 0.3 都差 0.1，并列选低灰度\\
1&0.7&2&到目标累计值 0.6 最近，差 0.1\\
2&0.8&2&到 0.6 和 1 都差 0.2，并列选低灰度\\
3&1&3&目标累计值恰好为 1\\
\bottomrule
\end{tabular}
\end{center}
得到映射 $0\to0,1\to2,2\to2,3\to3$，输出计数为 $2,0,6,2$，并不等于目标计数 $1,2,3,4$。原因仍是同灰度像素成组变换，且匹配位置有限。另一种常用规则是选择“第一个累计值不小于源累计值”的灰度；它可能得到不同答案，答题时应说明规则。

\textbf{学到这里就够用的认识：}匹配把两个 CDF 对应起来，不要求先真正保存一幅均衡图。理论上它与一维分布的单调传输有关，但一般离散图像的单值灰度表未必能精确实现任意目标分布。对入门计算，不必先学习最优传输理论。

\subsection{读直方图的原题 先描述 再选方法}

\subsubsection*{原题 B1 比较四幅图像的明暗和对比度}
\label{original:B1}

\textbf{题目。}比较题中 a、b、c、d 四个直方图对应图像的明暗与对比度，指出哪些需要处理，以及可选的方法。

\textbf{图源说明。}原讲义保留了题目和文字答案，但没有可核对的原始四幅题图。下面是为讲解重新绘制的\textbf{教学示意图}，字母只代表本页四种情况，不能用它断言缺失原图的精确形状或编号关系。

\begin{center}
\begin{tikzpicture}[x=.018cm,y=.47cm,>=Stealth,font=\small]
\begin{scope}
\draw[->] (0,0)--(275,0) node[right] {$r$};\draw[->] (0,0)--(0,5.2) node[above] {计数};
\foreach \k/\v in {20/1,35/2.7,50/4.3,65/3.5,80/1.5}{\draw[fill=black!15] (\k-7,0) rectangle (\k+7,\v);}
\node[below] at (0,0) {0};\node[below] at (255,0) {255};\node[above] at (135,5.1) {a 暗部集中};
\end{scope}
\begin{scope}[xshift=7.8cm]
\draw[->] (0,0)--(275,0) node[right] {$r$};\draw[->] (0,0)--(0,5.2) node[above] {计数};
\foreach \k/\v in {180/1,195/2.8,210/4.4,225/3.5,240/1.8}{\draw[fill=black!15] (\k-7,0) rectangle (\k+7,\v);}
\node[below] at (0,0) {0};\node[below] at (255,0) {255};\node[above] at (135,5.1) {b 亮部集中};
\end{scope}
\begin{scope}[yshift=-4.5cm]
\draw[->] (0,0)--(275,0) node[right] {$r$};\draw[->] (0,0)--(0,5.2) node[above] {计数};
\foreach \k/\v in {100/1.5,115/3.8,130/4.5,145/2.8,160/1}{\draw[fill=black!15] (\k-7,0) rectangle (\k+7,\v);}
\node[below] at (0,0) {0};\node[below] at (255,0) {255};\node[above] at (135,5.1) {c 中部窄范围};
\end{scope}
\begin{scope}[xshift=7.8cm,yshift=-4.5cm]
\draw[->] (0,0)--(275,0) node[right] {$r$};\draw[->] (0,0)--(0,5.2) node[above] {计数};
\foreach \k/\v in {15/1.3,40/2.1,65/2.5,90/2,115/2.8,140/3,165/2.4,190/2.6,215/2.1,240/1.4}{\draw[fill=black!15] (\k-9,0) rectangle (\k+9,\v);}
\node[below] at (0,0) {0};\node[below] at (255,0) {255};\node[above] at (135,5.1) {d 较宽范围};
\end{scope}
\end{tikzpicture}
\end{center}

\textbf{第一步，找横向位置。}a 的主要像素靠近黑端，偏暗；b 靠近白端，偏亮；c 在中部，明暗适中；d 涵盖暗、中、亮范围。

\textbf{第二步，看主要分布的宽度。}a、b、c 都较窄，通常整体对比度不足。d 跨度较宽，有较大的全局灰度范围，但这还不能保证每个局部都清楚。

\textbf{第三步，把方法与问题配对。}a 可以拉伸暗部范围，若还需提亮可试规范化 $\gamma<1$；b 可试 $\gamma>1$ 调暗并温和拉开未饱和的亮部差异；c 可用分段线性拉伸；d 不应仅因为有增强算法就强行处理，应先看实际图像和任务是否已经满意。

\textbf{答题边界。}对本页示意，a、b、c 更值得检查和处理。不能仅凭 d 的宽直方图选择“中值去噪”，因为直方图没有告诉我们它是否有椒盐噪声。缺失原始题图时，原题 a--d 的逐项结论仍需与原图核对。

\subsubsection*{原题 B2 夜间道路的双峰直方图}
\label{original:B2}

\textbf{题目。}夜间道路监控图像的直方图呈“两个高耸的峰，中间接近平坦”，说明什么？怎样增强？

\textbf{先解释数据。}两个峰说明至少有两组灰度范围贡献了较多像素。结合夜间道路场景，一个峰可能来自暗路面或天空，另一个可能来自车灯、反光牌等亮区域；这是根据场景作出的合理解释，不能只凭峰的个数确定物体种类。若题目的“平坦”指接近横轴的低谷，说明中间灰度较少；若只是高度不变，则还要看平台的具体高度。\textbf{不能据此认定细节已丢失}。

\textbf{再明确目标。}若只想检测亮灯位置，可以考虑分割亮区域；若想看清暗处行人，重点应是暗部可见性；若需要保持监控画面自然，则还要控制明暗关系和噪声。

\textbf{然后给出方案。}先检查暗区噪声与亮区饱和情况。若照明不均且暗区有可用细节，可试温和 CLAHE；若整体只是偏暗，可试规范化 $\gamma<1$ 或分段曲线，并检查车灯附近是否失去层次。需要去噪时，应按噪声特点选择方法，不能默认所有夜景都用中值滤波。

\textbf{最后说明风险。}全局 HE、CLAHE 和伽马都可能把原有噪声变得更明显。CLAHE 的限制机制降低过度增强风险，却不是“不放大噪声”的保证。完全饱和的车灯内部，也不能靠这些灰度变换恢复真实纹理。

\subsubsection*{原题 C1 根据暗人物与亮投影屏幕画直方图}
\label{original:C1}

\textbf{题目。}左侧是较暗的人物与背景，右侧是一块较亮的投影屏幕。画出直方图示意，并提出增强建议。

\textbf{图源说明。}原始题图未提供，因此无法知道峰的精确灰度、宽度、高度或面积。以下只根据文字场景新画示意，\textbf{不是原题图像的真实统计结果}。原文中的具体峰位数字缺乏可核对图像，本节不把它们当作事实。

\begin{center}
\begin{tikzpicture}[x=.042cm,y=.62cm,>=Stealth,font=\small]
\draw[->] (0,0)--(272,0) node[right] {灰度 $r$};
\draw[->] (0,0)--(0,5.2) node[above] {像素数};
\draw[thick,fill=black!12] (0,0)--(6,.2)--(12,1.4)--(20,2.9)--(30,3.6)--(40,3.1)--(52,1.6)--(66,.4)--(90,.15)--(130,.12)--(163,.15)--(182,.5)--(192,2.5)--(199,4.2)--(205,3.2)--(216,.5)--(230,.15)--(255,0)--cycle;
\node[align=center] at (60,4.5) {较暗人物与背景\\低灰度的一组像素};
\node[align=center] at (205,5) {较亮屏幕\\高灰度的一组像素};
\node[below] at (0,0) {0};\node[below] at (255,0) {255};
\end{tikzpicture}
\end{center}

\textbf{作图第 1 步：}先画灰度横轴和计数纵轴。横轴左黑右白；纵轴越高，说明该灰度的像素越多。

\textbf{作图第 2 步：}在较暗一侧画一组隆起，在较亮一侧画另一组隆起。屏幕若亮度比较接近，可能形成较窄峰；人物与背景的材质和照明不一，暗峰可能更宽。这些是条件性的解释，而非精确测量。

\textbf{作图第 3 步：}不要把“图像左侧的人物”理解成直方图必然左侧。人物的峰在左边，是因为\textbf{人物较暗}，而不是因为人物位于画面左边。直方图不保存空间位置。

\textbf{方法建议：}若目的是看清人物，可在暗部使用有保护亮部的分段曲线，或尝试温和 CLAHE；若还需读取屏幕文字，应同时检查亮区是否饱和，以及增强后文字笔画是否变差。对不同区域分开处理再平滑过渡也是一种方案。全局 HE 可以作比较基线，但不能在没有测试的情况下保证它或 CLAHE 一定更好。

\subsection{HE 的改进方法 先看它试图解决哪个问题}

学这一组名称时，不必背成“后一代全面替代前一代”的时间线。它们处理的问题并不完全相同：有的关心局部照明，有的减少平均灰度变化，有的约束过度增强，还有的借助光照模型或学习模型。下面只讲机制与边界，不用未经核对的年份帮助记忆。

\subsubsection*{全局 HE 常见的三个局限}

\begin{enumerate}[leftmargin=*]
\item \textbf{小区域可能被大背景主导。}若全图大部分是天空，按全图统计得到的曲线未必适合只占少量像素的车牌。
\item \textbf{噪声可能变得明显。}在平坦暗区，灰度 30、31、32 的差别可能主要是噪声；若它们映射到 40、100、170，噪声就被拉大。
\item \textbf{整体明暗可能改变。}HE 没有强制“输出平均灰度等于输入平均灰度”。想保留夜景感觉或照片自然感时，这种变化可能不合适。
\end{enumerate}

\subsubsection*{AHE 让每个局部按自己的统计来处理}

\textbf{自适应直方图均衡化（adaptive histogram equalization，AHE）}用像素周围的局部直方图来确定映射。最直接的理解是：移动一个窗口，统计窗口内的灰度，为当前位置计算均衡输出。高效实现也可以划分若干小块，再结合相邻区域的映射。

与全局 HE 相比，同为灰度 80 的像素，在暗角和亮处可以得到不同输出，因为周围的灰度分布不同。它因此能适应空间上的照明差异。

但局部几乎全是灰度 30 时，极少量的 29 和 31 可能只是噪声；AHE 仍可能把这个狭窄分布拉开。窗口越小，局部统计越容易被少量像素影响；窗口过大，又逐渐失去对小范围变化的适应能力。

\subsubsection*{CLAHE 给局部增强加上限制}

\textbf{限制对比度的自适应直方图均衡化（contrast-limited adaptive histogram equalization，CLAHE）}在局部均衡之前限制过高的直方图柱子，降低由尖峰引起的过强灰度拉伸。常见分块实现包含四步：
\begin{enumerate}[leftmargin=*]
\item \textbf{分块并计数。}把图像划分为小块，每块统计自己的直方图。
\item \textbf{裁剪并再分配。}设置\textbf{裁剪限值（clip limit）}，超过限值的计数先截下来，再按实现规则分配给各灰度级。截下的是用于构建映射的统计量，不是删除那些像素。
\item \textbf{累计并建映射。}对修改后的直方图求 CDF，为每个小块建立灰度对照表。
\item \textbf{在相邻映射间插值。}处理一个像素时，结合周围小块映射的结果，使过渡平滑。常用\textbf{双线性插值（bilinear interpolation）}，本质上按水平和竖直距离加权平均。
\end{enumerate}

\begin{center}
\begin{tikzpicture}[>=Stealth,font=\small]
\draw (0,0) rectangle (4.4,3.6);\draw (2.2,0)--(2.2,3.6);\draw (0,1.8)--(4.4,1.8);
\node[align=center] at (1.1,2.7) {块 A\\局部映射 $T_A$};
\node[align=center] at (3.3,2.7) {块 B\\局部映射 $T_B$};
\node[align=center] at (1.1,.9) {块 C\\局部映射 $T_C$};
\node[align=center] at (3.3,.9) {块 D\\局部映射 $T_D$};
\fill (2.25,1.85) circle (2pt);\node[above right,fill=white,inner sep=1pt] at (2.25,1.85) {$P$};
\node[align=center] at (2.2,-.5) {像素 $P$ 结合相邻块的映射\\不是给每块强行贴一张不同亮度的图};
\draw[->,thick] (4.8,1.8)--(5.6,1.8);
\draw[->] (6,0)--(11.2,0) node[right] {灰度};\draw[->] (6,0)--(6,3.7) node[above] {计数};
\foreach \x/\h in {6.5/.6,7.2/1.1,7.9/3.3,8.6/2.9,9.3/1.3,10/.7}{\draw[fill=black!12] (\x-.25,0) rectangle (\x+.25,\h);}
\draw[dashed,thick] (6,2)--(10.8,2);\node[right,align=left] at (10.8,2) {裁剪\\限值};
\draw[pattern=north east lines] (7.65,2) rectangle (8.15,3.3);
\draw[pattern=north east lines] (8.35,2) rectangle (8.85,2.9);
\node[align=center] at (8.6,-.5) {斜线部分重新分配\\再用修改后的计数求 CDF};
\end{tikzpicture}
\end{center}

\textbf{一遍裁剪的算术例子。}局部四个灰度的计数是 $[1,1,6,0]$，共 8 个。若限值为 3，先得到 $[1,1,3,0]$，截下 3 个计数。为演示整数再分配，假设把它们分别加给第 0、1、3 级，得到 $[2,2,3,1]$，总数仍为 8。实际程序可能按更细的规则分配，并在需要时再次处理超限，不能把这一个演示当成所有软件的唯一实现。

这张新表的累计比例是 $[2/8,4/8,7/8,8/8]$，再乘输出最大灰度并取整，就得到该块映射。\textbf{裁剪的是用于求映射的直方图，不意味着输出直方图每根柱子都低于限值。}

\textbf{插值也可以手算。}某像素用四个邻块映射分别会得到 80、100、120、140；若位于四个块中心构成区域的正中央，每个结果权重相同，输出便是 $(80+100+120+140)/4=110$。靠近某个块时，那个块的结果权重更大。图中仅示意四块与像素的关系，完整实现还需规定图像外边缘的处理方式。

\textbf{限制是什么。}CLAHE 可抑制 AHE 的过度噪声放大，并借助插值减轻块状接缝，但不等于去噪器，也不保证无噪声放大、无伪影或识别率上升。块尺寸和限值都要测试。不同软件对限值的单位、归一化方式不同，不能直接照搬另一个软件的参数数值。

\subsubsection*{BBHE 与 DSIHE 按灰度分组以减少亮度漂移}

先定义\textbf{平均灰度（mean gray level）}：把所有像素灰度相加，再除以像素总数，记为 $\mu$。已知直方图时，不必逐个像素加，可用
\[
\mu=\frac{0n_0+1n_1+\cdots+(L-1)n_{L-1}}{n}.
\]
例如前面的 10 像素图平均灰度为 $(0\times2+1\times5+2\times1+3\times2)/10=13/10=1.3$。

\textbf{亮度保持双直方图均衡化（brightness-preserving bi-histogram equalization，BBHE）}以输入均值为依据选分割灰度，把直方图分成较暗和较亮两组，分别在各自子范围内均衡。名称中的“亮度保持”表达设计目标，不是任意数字图像上严格相等的数学保证。

为看清机制，取分割阈值 $t=1$，低组为灰度 0、1，高组为 2、3。低组共有 $2+5=7$ 个像素，在低组内部重新算概率：$2/7,5/7$，不是仍除以 10。低组累计值是 $2/7,1$，映射到输出范围 $[0,1]$。高组共 3 个，累计值是 $1/3,1$，映射到 $[2,3]$。

若采用本节的四舍五入规则，低组输出为 $\operatorname{round}(2/7)=0$ 和 1；高组输出为 $\operatorname{round}(2+1/3)=2$ 和 3。这个特定例子恰好不变，但其他输入未必如此。

\textbf{一个直接反例。}设允许灰度为 0 到 7，输入四个像素为 $0,0,4,4$，均值为 2。按 $[0,2]$ 与 $[3,7]$ 分组，各组只有一个出现的灰度，其组内 CDF 为 1；用未减最小 CDF 的子区间均衡，0 映到 2，4 映到 7，输出为 $2,2,7,7$，均值变为 4.5。这个反例足以否定“BBHE 对所有离散图像精确保均值”。实现对单一灰度子组另有约定时结果可能不同，也更说明结论必须带条件。

\textbf{原讲义中那个均值条件从哪里来。}原文给出 $P(r\le m)=1-m/255$。这不是一般离散图像公式，但可以在特定理想模型中推导：设灰度连续地落在 $[0,255]$，输入均值为 $m$，按 $m$ 分成两组；低组比例为 $q=P(r\le m)$，高组比例便是 $1-q$。再假定处理后低组在 $[0,m]$ 完全均匀、高组在 $[m,255]$ 完全均匀，且两组均非空。

均匀区间的平均值是两端的平均，因此两组输出均值分别为 $m/2$ 和 $(m+255)/2$。全图均值要按组内像素比例加权：
\[
\mu_{\rm out}
=q\frac m2+(1-q)\frac{m+255}{2}
=\frac{m+255-255q}{2}.
\]
若要求它等于输入均值 $m$，两边乘 2 得 $m+255-255q=2m$，移项得 $255q=255-m$，再除以 255，才得到 $q=1-m/255$。所以该条件仅对上述\textbf{连续、子范围完全均匀}的理想模型成立。实际 BBHE 的整数阈值、上下子范围端点、取整和非均匀离散输出，都使这个推导不能直接当成通用定理。

\textbf{等面积双子图像直方图均衡化（dualistic sub-image histogram equalization，DSIHE）}改用中位数附近的灰度分割。\textbf{中位数（median）}是把像素值从小到大排列后处于中间的位置；从 CDF 看，就是累计比例到达约 0.5 的地方。分割后同样分别均衡，目标是让两组像素数尽量接近。

“按灰度分组”与 AHE 的“按空间分块”不同：BBHE 的暗组像素可能散布在整幅图中。另一个细节是：若某一灰度占比很大，不拆开同值像素就可能无法把像素数恰好二等分，不能把 DSIHE 的两组说成一定严格等大。

\subsubsection*{RMSHE 与 MMBEBHE 进一步控制平均亮度变化}

\textbf{递归均值分割直方图均衡化（recursive mean-separate histogram equalization，RMSHE）}把 BBHE 的分割继续做下去。先按全图均值分两组；然后在每组内部重新求均值，再各分两组。这个“对得到的子问题重复同样操作”的过程叫\textbf{递归（recursion）}。

若每组都能有效分裂，分 1 层最多有 2 组，分 2 层最多 4 组，分 $d$ 层最多 $2^d$ 组，$d$ 是分裂层数。最后每组只在自己的灰度范围内均衡。分组细化可限制灰度移动幅度，但也可能使增强变弱；数字图像还有空组、单值组和取整问题，不能宣称递归层数越多视觉效果必然越好。

\textbf{最小平均亮度误差双直方图均衡化（minimum mean brightness error bi-histogram equalization，MMBEBHE）}尝试用使亮度变化较小的分割阈值。先定义\textbf{绝对平均亮度误差（absolute mean brightness error，AMBE）}：输入均值记为 $\mu_{\rm in}$，输出均值记为 $\mu_{\rm out}$，则
\[
\operatorname{AMBE}=|\mu_{\rm out}-\mu_{\rm in}|.
\]
竖线表示绝对值。例如输入均值 80、输出均值 95，误差为 $|95-80|=15$；输出均值 70 时误差为 10，不是负 10。

MMBEBHE 的核心是比较候选分割阈值对应的亮度误差，选择较优阈值，再分别均衡。理解时可以想象“尝试几种切法，比较输出均值，选择误差最小的切法”；实际算法可用高效计算避免生成所有完整输出图。\textbf{最小只针对所比较的阈值、误差定义和算法模型，不代表图像各方面最好，也不保证误差为零。}一幅完全不处理的图 AMBE 就为零，所以这个指标不能单独衡量增强质量。

\subsubsection*{WTHE 先修改直方图的权重}

\textbf{加权阈值直方图均衡化（weighted thresholded histogram equalization，WTHE）}先对直方图概率设置阈值并作加权，再用修改后的分布计算映射。它试图避免少数很高的柱子过度控制映射，也控制低频数灰度对增强的影响。

这里的“加权”是改变各灰度在构建映射时的贡献，不是把图像某些像素变重。修改后通常还要重新归一化，使用于累计的概率总和为 1。它与 CLAHE 都可能限制直方图的影响，但\textbf{WTHE 本身不等于空间分块}；具体权重和阈值需看所采用的版本。入门时先记“先改分布，再求 CDF”，不必背没有说明参数含义的公式。

\subsubsection*{Retinex 与 LIME 从光照和物体反射来理解暗图}

\textbf{视网膜皮层理论及其图像增强方法（Retinex）}常用一个简化模型理解观察到的图像：同一个物体，照得亮就更亮，照得暗就更暗；图像的明暗同时受照明和物体表面反射特性影响。

把位置记为 $(x,y)$，观察强度记为 $I(x,y)$，照明分量记为 $\ell(x,y)$，反射分量记为 $R(x,y)$，便有模型
\[
I(x,y)=\ell(x,y)R(x,y).
\]
这里的 $\ell$ 是小写字母 ell，表示照明，不是前文灰度级数 $L$。相乘表示照明与反射共同决定强度；例如反射分量为 0.5，照明从 0.2 变为 0.8，观察强度便从 0.1 变为 0.4。

Retinex 类方法估计或消除部分照明影响，以改善可见性和颜色表现。常见方法借助局部平滑估计背景照明，并在对数域中处理乘积关系。但单幅图像通常不能唯一确定照明和反射：$0.2=0.4\times0.5=0.5\times0.4$。因此必须引入平滑、边缘或颜色等假设，且可能出现光晕、噪声增强或颜色失真。不能把“物理分解”误解成自动获得真实光照。

\textbf{基于照明图估计的低照度增强（low-light image enhancement via illumination map estimation，LIME）}以照明图估计为重点。它可从彩色图每个像素的红、绿、蓝三个通道中的最大值建立初始照明估计，再利用空间结构对照明图进行细化，根据估计结果增强暗处。所谓\textbf{通道（channel）}，就是分别记录某种颜色分量的一张数值图。

直观上，同样的观察灰度若处于明显较暗的照明区域，就需要不同于亮区的调整。LIME 的光照估计与空间结构有关，并不是一条全图共用的普通伽马曲线。除以很小的照明估计时，原有噪声也可能被放大，因此照明正则化、数值保护及必要的去噪都很重要。

\subsubsection*{Zero-DCE 让模型预测局部曲线参数}

\textbf{零参考深度曲线估计（zero-reference deep curve estimation，Zero-DCE）}用神经网络预测增强曲线的参数。\textbf{神经网络（neural network）}在这里可先理解为一个通过训练学会“根据图像给出参数”的函数，不需要先掌握其全部内部结构。

其基础增强曲线可写成
\[
v=u+\alpha u(1-u).
\]
$u$ 是 0 到 1 的归一化输入，$v$ 是输出，$\alpha$ 是控制弯曲程度的参数。先算 $1-u$，再与 $u$ 相乘，乘以 $\alpha$，最后加回 $u$。当 $u=0.2,\alpha=0.5$ 时，$v=0.2+0.5\times0.2\times0.8=0.28$，于是变亮。若 $\alpha=-0.5$，则得到 0.12，变暗。当 $-1\le\alpha\le1$ 时，曲线在 $0\le u\le1$ 内保持输出范围并且单调不减，即输入较大时输出不会较小；多次重复曲线调整可产生更灵活的变换。

关键区别是：参数可以随位置、通道及迭代步骤变化，由网络结合图像内容预测。因此最终曲线作用于像素，不代表整个算法只看该像素自己的输入；\textbf{它不是一条固定的全局 $T(r)$}。

“零参考”表示训练不需要正常曝光参考图作为监督目标，并不表示“不需要训练数据”“不需要训练”或“任何暗图都能无损恢复”。模型仍需要图像和约束增强效果的损失函数；噪声、极暗细节和饱和信息仍有困难。

\subsubsection*{用一张对照表整理方法}
\begin{center}
\small
\begin{tabularx}{\linewidth}{@{}p{2.4cm}p{4.4cm}X@{}}
\toprule
方法 & 主要改变什么 & 先记住的限制\\
\midrule
HE & 全图共用累计映射 & 可改均值、放大噪声；离散输出不必均匀\\
AHE & 根据局部统计选映射 & 平坦局部易过度增强\\
CLAHE & 局部直方图限幅，再平滑结合映射 & 仍需调参，不能保证去噪\\
BBHE / DSIHE & 按均值 / 中位数分灰度组 & 不同于空间分块，不保证精确保均值\\
RMSHE & 反复按各组均值细分 & 细分更多可能减弱增强\\
MMBEBHE & 按亮度误差选分割阈值 & 均值误差小不等于整体质量高\\
WTHE & 修改直方图权重与阈值 & 具体参数影响增强和细节\\
Retinex、LIME & 估计或调整照明的影响 & 分解依赖假设，可能有光晕与噪声\\
Zero-DCE & 网络预测空间变化的曲线参数 & 仍有训练、数据与泛化要求\\
\bottomrule
\end{tabularx}
\end{center}

\subsection{原题 G3 低对比度车牌在 HE CLAHE 和伽马中怎么选}
\label{original:G3}

\textbf{题目。}对一张低对比度车牌图像，要在 HE、CLAHE、伽马变换中选一种，服务于后续字符识别。你选哪个，为什么？

\textbf{第 1 步，确认目前不知道什么。}“低对比度”本身不足以唯一决定方法，还需知道它是整体偏暗、局部照明不均、噪声重，还是已经模糊或饱和。题目不给图像时，要把选择的前提写出来，不能假装知道车牌实际照明。

\textbf{第 2 步，给出有条件的主方案。}若车牌受到明显不均匀照明，可把温和 CLAHE 作为优先测试方案。它根据局部统计增强字符与背景的差异，并限制过强拉伸；相邻块映射的插值有助于减轻分块接缝。

\textbf{第 3 步，与其他候选公平比较。}若主要问题是整体偏暗且照明均匀，规范化 $\gamma<1$ 可能已经足够，且参数简单。HE 可作为自动全局对照方案；它不利用位置，可能被大背景主导，但也不能断言一定失败。伽马能改变某些灰度区间的对比度，只是不具有 CLAHE 的空间自适应性。

\textbf{第 4 步，写出真正的验收方式。}尽量在车牌区域比较原图、HE、CLAHE 和伽马结果；保持后续识别器和其他处理条件一致，用有正确答案的样例检查字符或整牌识别正确率，同时查看笔画断裂、粘连、噪声与饱和。不能用“看起来更亮”代替“识别更准确”。

\textbf{可写在答卷上的答案：}若低对比度主要来自局部照明不均，我优先测试 CLAHE，因为它具有局部适应性和限制过度增强的机制；采用适当块尺寸与裁剪限值，并以识别结果验证。若只是均匀的整体偏暗，伽马变换也可能更合适。最终方法应由图像条件和字符识别效果决定，CLAHE 不保证始终最优。

\subsection{本节自查 把公式变成能解释的步骤}

不看前文，试着回答以下问题，再用后面的提示检查。
\begin{enumerate}[leftmargin=*]
\item 8 位图中 $L$ 为什么是 256，而最大灰度是 255？\par
\textbf{检查：}共有 256 个整数，从 0 开始数到 255。
\item 为何 $\gamma=1/2$ 有时算得 8，有时算得 128？\par
\textbf{检查：}对 64 直接开方得 8；对 $64/255$ 开方再乘 255，才是本节保持范围的规范化伽马，约得 128。
\item 灰度 2 的概率与累计概率有何不同？\par
\textbf{检查：}前者只数灰度 2，后者把 0、1、2 都计入，再除以总数。
\item E1 中输出 6 为什么有 985 个像素？\par
\textbf{检查：}输入 3 与 4 都映射成 6，所以把其计数 $656+329$ 相加。
\item 两张图直方图相同，为什么仍可能完全不像？\par
\textbf{检查：}直方图没有保存像素位置和排列结构。
\item CLAHE 比 AHE 多一个限制机制，为什么仍要检查噪声？\par
\textbf{检查：}限制增强幅度只是减少过度放大风险，不等于估计并去掉噪声。
\end{enumerate}
